\subsection{R-28: higher-variable determinant dimensions and common free points}

\paragraph{Statement R-28.}
For every $n\geq3$ and $r\geq4$,
\begin{equation}\label{R28:det_dimension}
 \dim\mathcal D_{n,r}=(r-2)n^2+2.
\end{equation}
For $m\geq17$, $r\geq4$, and $n>m$, the padded permanent
image has dimension
\begin{equation}\label{R28:pad_dimension}
 \dim\{\ell^{n-m}q:q\in\mathcal P_{m,r}\}^{\rm cl}
       =rm^2-2m+r+1.
\end{equation}

In the setting of G-22 with $n>m\geq r\geq3$ and
$\nu_r>0$, $\nu_1>(n-m)s$, its two enhanced cones have a
common stable point $w$ with the following stronger property:
\begin{equation}\label{R28:no_scaling}
 g\cdot w=cw,\quad g\in\mathrm{SL}_r,\quad c\in\mathbb C^*
       \quad\Longrightarrow\quad c=1.
\end{equation}
In fact the underlying padded form can be chosen with trivial
projective stabilizer in $\mathrm{PGL}_r$.
If homogeneous generators of $\mathbb C[W]^{\mathrm{SL}_r}$
have positive degrees at most an integer $\beta$, those which
are nonzero at this common point have degrees of greatest common
divisor one. Consequently both multiplicities $m_{j\nu}$ are
positive for every integer $j\geq\beta^2$.
This includes the endpoint $r=m$ and padding exponent one.

\paragraph{Background used.}
Openness of the stable locus and separation of distinct closed
orbits are cited from Mukai (2003), Proposition~5.15, printed
p.~166, and Theorem~5.3 with Corollary~5.4, printed pp.~161--162
(G-B15). The common stable product point and exact invariant
extraction are G-22. The cyclic perturbation, the explicit
irreducibility and stabilizer witnesses, and the extension of
the image dimension are proved here.

\begin{proof}
We work over $\mathbb C$ and treat the dimension and stability statements in turn.

\paragraph{1. Exact determinant dimension after adding variables.}
R-17 gives $\dim\mathcal D_{n,4}=2n^2+2$. The cyclic
cofactor independence established in R-23 also applies to
determinants: each complementary matching at that cyclic
pencil is unique, so its determinant cofactor is the same
polynomial as its permanent cofactor up to a nonzero sign.
Thus the $n^2$ determinant cofactors are independent on a
nonempty open subset of the quaternary parameter space.
Intersect this with the nonempty open where the differential
has rank $2n^2+2$, and choose one pencil $A$ in the intersection.

Regard $A$ as an $r$-variable pencil with zero coefficients
on the additional variables. For each new variable $x_k$,
varying its $n^2$ matrix coefficients gives $x_k$ times the
independent cofactor family. These blocks have dimension
$n^2$ each. Their multidegrees in the new variables are
different from one another and from the old block. The
differential therefore has rank
\[
 2n^2+2+(r-4)n^2=(r-2)n^2+2.
\]
For the opposite inequality, the left--right action of
$\mathrm{SL}_n\times\mathrm{SL}_n$ preserves the determinant
and generically has finite stabilizer on pencils with at
least three coefficient matrices. For example, take one
coefficient to be $I$, a second diagonal with distinct entries,
and a third with every entry nonzero. The first forces the
two acting matrices to be inverses, the second forces the
remaining one to be diagonal, and the third forces it to be
scalar. The special-linear condition then leaves only a
finite central group. The generic fiber thus contains an
orbit of dimension $2n^2-2$, giving the upper bound
$rn^2-(2n^2-2)$. This proves \eqref{R28:det_dimension}.
The argument is valid even when $r>n^2$; it concerns a
coefficient map and does not require that its variables be
obtained by a quotient by linear equations.

The padding argument of R-14 works in $r$ variables. A general
member of $\mathcal P_{m,r}$ is irreducible: this cone contains
the Eisenstein polynomial $x_1^m+x_2^{m-1}x_3$, supplied by
the diagonal-plus-cycle pencil. Irreducibility is a nonempty
open condition on degree-$m$ forms. The projective multiplication
map from $\mathbb P^{r-1}\times\mathbb P(\mathcal P_{m,r})$
is generically injective, since the nonlinear irreducible
factor and the padding line can be recovered separately from
their product. Its affine image therefore has dimension
$\dim\mathcal P_{m,r}+r-1$. Substitute G-21 to obtain
\eqref{R28:pad_dimension}.

\paragraph{2. One common cyclic family has all three needed open sets.}
Put $a=n-m$, retain the fixed extremal vector $v_\eta$ of
G-22, and consider the irreducible affine parameter space of
linear forms $L_1,\ldots,L_m,h,\ell$ and a scalar $\varepsilon$.
Set
\begin{equation}\label{R28:cyclic_family}
 q=\prod_{j=1}^mL_j+\varepsilon h^m,\qquad
 F=\ell^a q,\qquad w=F^{\otimes s}\otimes v_\eta.
\end{equation}
Every $q$ belongs to both the size-$m$ permanent and determinant
images. Indeed, the matrix with diagonal $L_j$, entries $h$
immediately above the diagonal, and corner entry $\varepsilon h$
at $(m,1)$ has exactly two perfect matchings. Its permanent
is $q$; changing that corner to $(-1)^{m-1}\varepsilon h$
gives determinant $q$. Adding $a$ diagonal entries $\ell$
then gives the required degree-$n$ determinant. Thus every
$w$ in this family lies in both enhanced cones.

The inverse image of the stable locus in $W$ is nonempty and
open in this parameter space. It is nonempty at
$\varepsilon=0$ by the Fourier-frame construction in G-22,
using its $L_j$, padding line, and the same fixed $v_\eta$.
We will intersect it with the following two additional
nonempty open conditions on $q$:
\[
 q\text{ is irreducible},\qquad
 \operatorname{Stab}_{\mathrm{PGL}_r}([q])\text{ is finite}.
\]
The first is open because reducible degree-$m$ forms are a
finite union of closed projective multiplication images.
The second is open by the maximal-rank condition on the
coefficient matrix of the infinitesimal equation
$Bq=cq$: its only solutions should be scalar matrices $B$.
Over characteristic zero, zero projective stabilizer Lie
algebra means a finite algebraic stabilizer.

Here is one explicit witness for both conditions in every
$m\geq r\geq3$. In coordinates $x_1,\ldots,x_r$, put
\begin{equation}\label{R28:explicit_q}
 P=x_1x_2\cdots x_{r-1}x_r^{m-r+1},\qquad
 z=x_1+\cdots+x_r,\qquad q_*=P+\varepsilon z^m,
 \quad\varepsilon\ne0.
\end{equation}
This is a specialization of \eqref{R28:cyclic_family}, allowing
repeated diagonal factors. Regard $z,x_1,\ldots,x_{r-1}$ as
independent coordinates. Then
\[
 q_* = \varepsilon z^m+
       (x_1\cdots x_{r-1})
          (z-x_1-\cdots-x_{r-1})^{m-r+1}.
\]
As a polynomial in $z$ over
$\mathbb C[x_1,\ldots,x_{r-1}]$, its leading coefficient is
a unit, every other coefficient is divisible by $x_1$, and
its constant coefficient is divisible by $x_1$ exactly once.
The latter assertion uses $r\geq3$, so
$x_1+\cdots+x_{r-1}$ is not divisible by $x_1$.
Eisenstein at $x_1$ proves irreducibility.

For the stabilizer, suppose a linear infinitesimal operator
$B$ satisfies $Bq_*=cq_*$. No pure monomial $x_i^m$ occurs
in $BP$: all $r\geq3$ variables occur in $P$, and one
infinitesimal replacement can remove at most one of them.
Writing $Bz=\sum b_i x_i$, comparison of the pure monomials
in $\varepsilon m z^{m-1}Bz=c\varepsilon z^m$ gives
$b_i=c/m$ for every $i$. Hence $Bz=(c/m)z$, and the remaining
equation is $BP=cP$. Each off-diagonal matrix entry contributes
a different monomial $P x_i/x_j$, with nonzero coefficient
because every exponent of $P$ is positive. It follows that
all off-diagonal entries vanish. The equation for $Bz$ then
makes every diagonal entry equal to $c/m$. Thus $B$ is scalar,
as required.

The three nonempty open subsets just constructed intersect
because the parameter space is irreducible. Choose a tuple in
their intersection and fix its irreducible $q$ and its auxiliary
vector. Its projective automorphism group $H$ is finite.

\paragraph{3. Varying the padding line removes every projective symmetry.}
At the tuple just chosen, the line $\ell$ belongs to a
nonempty open set of padding lines for which
$(\ell^a q)^{\otimes s}\otimes v_\eta$ is stable.
Every nonidentity element of $H\subset\mathrm{PGL}_r$ has
a proper closed fixed locus on the projective space of linear
forms. Avoid the finite union of these loci while remaining
in that stable open. For this choice of $\ell$, the projective
stabilizer of $F=\ell^a q$ is trivial. Indeed, unique
factorization forces a transformation preserving $[F]$ to
preserve the nonlinear irreducible factor $[q]$ and the linear
factor $[\ell]$ separately. Thus it belongs to $H$ and fixes
$[\ell]$, which forces the identity. This argument does not
require $a\geq2$: the two factors are distinguished by their
degrees even when $a=1$.

If $g\in\mathrm{SL}_r$ satisfies $gw=cw$, the two nonzero
decomposable tensor factors show that $g$ preserves the line
of $F^{\otimes s}$ and hence the line of $F$. The Veronese
map is injective on lines. Therefore $g$ is central, say
$g=tI$ with $t^r=1$. Its action on $W$ is multiplication by
\[
 t^{ns-e}=t^{r\nu_r}=1.
\]
This proves \eqref{R28:no_scaling}. In particular $w$ and
every nonzero scalar multiple of $w$ have closed orbits,
and two different scalar multiples cannot lie in the same
$\mathrm{SL}_r$ orbit.

\paragraph{4. Common invariant degrees and occurrence at every large stretch.}
Let $f_1,\ldots,f_t$ be homogeneous positive-degree generators
of the ambient invariant ring, with degrees at most $\beta$,
and keep the indices for which $f_i(w)\ne0$. There is at least
one: otherwise the closed orbits of $w$ and zero would not
be separated by invariants. If the kept degrees had greatest
common divisor $d>1$, a primitive $d$th root $c\ne1$ would
satisfy $f_i(cw)=f_i(w)$ for every generator, including those
vanishing at $w$. Every invariant would then agree at $w,cw$.
Closed-orbit separation would put them in the same orbit,
contradicting \eqref{R28:no_scaling}. The greatest common
divisor is therefore one.

The elementary numerical-semigroup bound proved in G-18
now applies: positive generators bounded by $\beta$ and of
greatest common divisor one generate every integer at least
$\beta^2$. Products of the kept invariant generators are
nonzero at $w$, and hence nonzero on both enhanced cones.
G-22's exact invariant extraction identifies these degree-$j$
elements with the respective occurrences of $j\nu$ in original
coordinate degree $sj$. This proves the final assertion.
\end{proof}

\paragraph{Scope.}
The first two image dimensions are unconditional in their
stated ranges. The common free point requires the strict
direction hypotheses and $3\leq r\leq m$, as stated.
No conclusion is drawn here about the padding boundary
$\nu_1=(n-m)s$, where a remaining $(r-1)$-variable image can
already have nontrivial multiplicity differences. The result
supplies the absence of an invariant-degree period for the
higher-row comparison, not a positive extraction algorithm
for individual coefficients.
